\documentclass[conference]{IEEEtran}
\usepackage{cite}
\usepackage{amsmath,amssymb,amsfonts}
\usepackage{algorithmic}
\usepackage{graphicx}
\usepackage{textcomp}
\usepackage{xcolor}
\usepackage{url}

\begin{document}

\title{Apple Leaf Disease Detection System Using Multiple Machine Learning Approaches}

\author{\IEEEauthorblockN{Computer Vision Research Team}
\IEEEauthorblockA{Department of Computer Science\\
University\\
Email: author@university.edu}}

\maketitle

\begin{abstract}
This paper presents a comprehensive apple leaf disease detection system utilizing both deep learning and traditional machine learning approaches. The system employs multiple convolutional neural network architectures including MobileNetV2, ResNet50, DenseNet121, and EfficientNetB0, alongside classical machine learning models such as Support Vector Machine, Random Forest, and K-Nearest Neighbors. Additionally, system incorporates disease severity estimation using image processing techniques to quantify extent of infection on apple leaves. The proposed system achieves high accuracy rates across multiple disease categories including Apple Scab, Black Rot, Cedar Apple Rust, and Healthy leaves.
\end{abstract}

\begin{IEEEkeywords}
Apple disease detection, deep learning, machine learning, computer vision, image processing, severity estimation, convolutional neural networks, support vector machine
\end{IEEEkeywords}

\section{Introduction}
Apple cultivation faces significant challenges from various diseases that affect leaf health and crop yield. Early detection and accurate classification of these diseases are crucial for effective disease management and minimizing economic losses. Traditional methods relying on visual inspection by experts are time-consuming and subject to human error.

This research addresses these challenges by developing an automated disease detection system that utilizes multiple state-of-the-art deep learning architectures, incorporates traditional machine learning models with feature extraction, provides disease severity estimation using image processing, offers comprehensive model comparison and analysis, and delivers real-time predictions through an interactive web interface.

\section{Methodology}
\subsection{Dataset Description}
The dataset consists of apple leaf images categorized into four classes: Apple Scab (1,368 images), Black Rot (1,368 images), Cedar Apple Rust (1,200 images), and Healthy (1,404 images). The total dataset contains 5,340 training images and 2,222 validation images. All images are resized to 224x224 pixels and normalized to the range [0, 1].

\subsection{Proposed System}
The system consists of three main components: Deep Learning Models with multiple CNN architectures, Machine Learning Models with classical ML and feature extraction, and Severity Estimation using image processing for infection quantification.

The deep learning component includes MobileNetV2 with depthwise separable convolutions, ResNet50 with residual network architecture, DenseNet121 with dense convolutional network, and EfficientNetB0 with compound scaling. All models use pre-trained ImageNet weights with custom classification layers.

The machine learning component utilizes MobileNetV2 as a feature extractor, with Support Vector Machine using RBF kernel, Random Forest with 100 decision trees, and K-Nearest Neighbors with k=5 parameters.

The severity estimation component converts RGB images to HSV color space, detects infected regions using brown and yellow HSV ranges, creates masks for infected and leaf pixels, calculates infection ratio, and classifies severity levels as Mild (0-10%), Moderate (10-30%), Severe (30-60%), or Critical (60-100%).

\section{Results and Discussion}
\subsection{Model Performance}
The system achieved high accuracy across all tested models. Support Vector Machine achieved the highest accuracy of 99.06%, followed by EfficientNetB0 at 97.00%, KNN at 97.00%, DenseNet121 at 96.00%, Random Forest at 95.41%, ResNet50 at 95.00%, MobileNetV2 at 94.00%, and Custom CNN at 91.00%.

\subsection{Severity Analysis}
The disease severity estimation successfully quantifies infection levels using HSV-based image processing. The system provides visual indicators with color-coded severity levels: green for Mild, yellow for Moderate, orange for Severe, and red for Critical. The severity percentage is calculated as the ratio between infected pixels and total leaf pixels.

\subsection{System Evaluation}
The multi-model approach ensures robust predictions with consistent performance above 90% accuracy across all models. The web interface provides real-time predictions with comprehensive results display, model performance comparison, and statistical analysis. The severity estimation adds valuable information for disease management decisions.

\section{Conclusion}
This research presents a comprehensive apple leaf disease detection system that successfully combines deep learning and traditional machine learning approaches. The key contributions include multi-model approach with eight different models for robust predictions, high accuracy with best model achieving 99.06% accuracy, effective severity analysis with HSV-based infection quantification, real-time deployment with interactive web interface, and comprehensive analysis with model comparison and statistical insights.

Future work includes expanding dataset to include more apple disease varieties, implementing real-time video processing capabilities, developing mobile application for field use, integrating weather data for disease prediction, and exploring ensemble methods for improved accuracy.

\begin{thebibliography}{1}
\bibliographystyle{IEEEtran}

\end{document}
